% Basierend auf dem JMLR Template jmlr2e.
\documentclass[twoside,11pt,abbrvbib]{article}

\usepackage{jmlr2e}
\usepackage{amsmath}
\usepackage{parskip}
\usepackage[utf8]{inputenc}
\usepackage[T1]{fontenc}
\usepackage[ngerman]{babel}
%
% Definitions of handy macros can go here

\newcommand{\dataset}{{\cal D}}
\newcommand{\fracpartial}[2]{\frac{\partial #1}{\partial  #2}}

% Heading arguments are {volume}{year}{pages}{submitted}{published}{author-full-names}
\jmlrheading{}{}{}{}{}{} % Ignore this...




% Short headings should be running head and authors last names
\ShortHeadings{Abgekürzter Titel}{Mustermann}

\begin{document}

\title{Titel}

\author{\name Melinda Mustermann\\
    \addr Technische Universität Dortmund\\
    \email melinda.mustermann@tu-dortmund.de \\
}


\maketitle

\begin{abstract}%
Dies ist eine kurze Zusammenfassung der Inhalte des in deutscher Sprache
verfassten Dokuments. Die Zusammenfassung sollte 10-20 Zeilen umfassen.
\end{abstract}


\section{Einleitendes Kapitel}
Lorem ipsum dolor sit amet, consectetur adipiscing elit. Proin id enim eu dui pellentesque malesuada nec ut metus. Phasellus augue justo, iaculis a velit a, elementum elementum mi. Aliquam vel condimentum leo. Phasellus sollicitudin aliquet justo, eu venenatis odio scelerisque et. Vivamus lorem diam, sodales et facilisis et, placerat vitae enim. Aliquam quis ultricies leo. Praesent congue orci nisi, vitae sodales mi sagittis eget. Phasellus ex tellus, convallis in ex quis, iaculis vehicula metus. Pellentesque vitae tristique enim. Quisque consectetur odio justo, eget lobortis est tincidunt et. Sed bibendum elit sed felis scelerisque, id sagittis lacus ultricies. Vivamus molestie elit sapien, sit amet placerat enim semper sed. Curabitur auctor vulputate nulla in maximus. Pellentesque pretium vel metus et rutrum. Nunc ante neque, ultricies non dui vitae, aliquet viverra leo. Donec pulvinar, nisl vitae sagittis porta, turpis magna vehicula lorem, sed facilisis dui sapien at dui.
\section{Eingabe mit mathematischer Notation}
In diesem Abschnitt ist zu sehen, was mit Formeln \cite{chow:68} passiert
\begin{equation}
\lim x =  \theta + 152383.52,
\end{equation}
die mit der equation-Umgebung gesetzt werden.
\section{Lorem Ipsum}


Praesent malesuada erat sed massa sodales tincidunt. Morbi id pharetra quam. Cras magna leo, mollis nec velit ut, volutpat gravida erat. Nullam aliquam faucibus nisi ut commodo. In blandit eleifend mauris lobortis aliquam. Phasellus vehicula id sem eu dignissim. Integer convallis quis ante sed aliquet. Donec varius commodo eros, et lobortis sem dignissim ut.

Orci varius natoque penatibus et magnis dis parturient montes, nascetur ridiculus mus. Suspendisse vitae lacus ex. Etiam viverra mollis odio, vel tempus felis faucibus id. Curabitur ac lectus lectus. Mauris libero quam, eleifend porta leo nec, cursus euismod nunc. Phasellus consectetur quam vel odio eleifend tempus. Aliquam vitae pulvinar ligula. Etiam nec quam sit amet massa venenatis suscipit ac ac ante.

In hac habitasse platea dictumst. Suspendisse potenti. Proin sit amet elementum metus. Cras sollicitudin congue ultricies. Mauris tincidunt malesuada scelerisque. Maecenas vestibulum nec nibh non placerat. Donec in placerat lectus, quis volutpat risus. Aliquam mattis, dui eu blandit auctor, diam velit volutpat lorem, ut gravida neque magna vitae diam.

In a lorem vitae tellus pellentesque ullamcorper id nec tortor. Suspendisse lorem risus, venenatis nec quam eu, suscipit ullamcorper sem. Curabitur pellentesque leo nunc, a semper sem pulvinar nec. Suspendisse nunc dui, auctor ac vestibulum ut, luctus a sapien. Suspendisse sem sapien, venenatis eget urna ac, lobortis scelerisque nunc. Sed libero libero, molestie et enim ut, malesuada porttitor nisi. Lorem ipsum dolor sit amet, consectetur adipiscing elit. Praesent tortor eros, lacinia non metus ut, porttitor suscipit neque. In varius, dolor eu pretium sagittis, diam leo semper lorem, in feugiat neque lectus eu sapien. Proin sodales scelerisque nisl a molestie. Nunc varius convallis bibendum.

\bibliographystyle{plain}
\bibliography{sample}

\newpage

\appendix
\section*{Appendix A.}
\label{app:theorem}

Falls ihr im Anhang noch Informationen aufführen wollt.


\vskip 0.2in


\end{document}